\begin{textsample}{POS Dim 3 – sp – Score 14.00 – sp/sp\_inf\_aspectos\_pedagogicos.txt}  \label{ex:f3_pos_206}
ASPECTOS PEDAGÓGICOS Ambientes : tempos, espaços e materiais Na \textbf{instituição} de Educação \textbf{Infantil}, a \textbf{rotina} deve ser permeada por marcos que possam proporcionar à \textbf{criança} regularidade das ações, de modo a criar segurança, \textbf{conforto} e noção de organização temporal.
Desde o momento de acolhida até a despedida, o \textbf{dia} a \textbf{dia} do \textbf{bebê}, das \textbf{crianças} bem \textbf{pequenas} e das \textbf{crianças} \textbf{pequenas} na \textbf{Instituição} de Educação \textbf{Infantil}, é permeado de situações relacionadas ao atendimento de suas necessidades como : alimentação, \textbf{higiene}, descanso e de momentos com as propostas para o trabalho com os objetivos de aprendizagem e desenvolvimento.
Dentre essas situações, todas permanentes e carregadas de intencionalidade, estão a contação de histórias, as \textbf{brincadeiras} na área externa, os jogos simbólicos, entre outros.
Ao se garantir regularidade, as \textbf{crianças} vão atribuindo significados a estes momentos, tornando -os marcos da \textbf{rotina}.
As \textbf{crianças} que frequentam a escola em período integral, por exemplo, logo que chegam exploram o solário ou área externa ; ao retornar para a sala de referência, sempre encontram uma novidade trazida pelos professores ; podem, ainda, escolher entre descobrir o que há de novidade ou explorar o espaço da sala, que deve estar organizado de forma a considerar a autonomia, os interesses e as necessidades das \textbf{crianças} e as especificidades da faixa \textbf{etária}.
Após este momento de descoberta, inicia -se a preparação para a alimentação e as \textbf{crianças} sabem que em \textbf{pequenos} grupos serão acompanhadas ao refeitório.
Nesse caso, os professores precisam estar sensíveis àqueles que \textbf{demonstram} necessidade primeira de se alimentarem, seja por desinteresse nas atividades propostas, seja por mostrarem -se fatigados, com sono ou com fome.
Então, terão, após a alimentação, o próximo momento, o descanso e assim, sucessivamente.
Até a despedida, a \textbf{rotina} acontece de forma regular, permeada de intencionalidade educativa que é revelada na efetivação do planejamento de tempos e espaços dos professores.
É importante destacar que a organização dos tempos e espaços deve estar a favor dos \textbf{bebês} e das \textbf{crianças}, não sendo necessário que se ajustem forçosamente às demandas da \textbf{instituição}.
Além disso, os tempos de transição entre uma atividade e outra também devem ser planejados, de forma que os \textbf{bebês} e as \textbf{crianças} não tenham que ficar em longo tempo de espera.
Também é imprescindível ter clareza de que os \textbf{cuidados} nesta fase são necessidades intrínsecas ao educar e que trocas e banhos acontecem ao longo da \textbf{rotina} sempre que necessários, sem horas engessadas e demarcadas.
O cotidiano precisa estar explicitamente a favor das necessidades das \textbf{crianças}.
Organizar tempos e espaços voltados às necessidades e interesses das \textbf{crianças} é fundamental para se garantir uma educação construída que considere a \textbf{criança} como competente e curiosa.
Essa educação é construída por meio de uma \textbf{rotina} que valida a participação da \textbf{criança} nas mais diversas situações vivenciadas na escola, desde a acolhida até a despedida.
Nesse sentido, a \textbf{escuta} da \textbf{criança} em suas múltiplas linguagens se faz primordial para que de fato ela se sinta parte ativa na \textbf{instituição}.
A disposição de mobílias e materiais pelo espaço tem de ser um convite à exploração e à descoberta.
Por isso, privilegiar espaços de participação nas \textbf{brincadeiras} e nas tomadas de decisões são princípios que regem uma educação voltada aos seus interesses.

% matched lemmas: bebê, brincadeira, conforto, criança, cuidado, demonstrar, dia, escuta, etário, higiene, infantil, instituição, pequeno, rotina
\end{textsample}
